\documentclass[12pt,a4paper]{article}
\usepackage[utf8]{inputenc}
\usepackage[T1]{fontenc}
\usepackage{amsmath,amssymb,amsthm}
\usepackage{physics}
\usepackage{geometry}
\usepackage{enumitem}
\usepackage{hyperref}
\usepackage{fancyhdr}
\usepackage{titlesec}

\geometry{margin=1in}
\pagestyle{fancy}
\fancyhf{}
\rhead{Lecture 6}
\lhead{Quantum Mechanics --- The Theoretical Minimum}
\cfoot{\thepage}

\titleformat{\section}{\large\bfseries}{}{0em}{}
\titleformat{\subsection}{\normalsize\bfseries}{}{0em}{}

\newtheorem{definition}{Definition}
\newtheorem{theorem}{Theorem}

\title{\textbf{Lecture 6: Entanglement and Composite Systems} \\ \large Tensor Products, Two Spins, and the Collapse of the Wave Function}
\author{Based on Leonard Susskind's \textit{The Theoretical Minimum}}
\date{}

\begin{document}
\maketitle
\thispagestyle{fancy}

% ============================================================
\section{1. Photon Emission and Energy Levels}
% ============================================================

Consider a spin with Hamiltonian $H=\frac{\omega}{2}\sigma_z$, giving energy levels $E_\pm = \pm\frac{\omega}{2}$.

\noindent\textbf{Classical picture:} A classical spinning magnetic dipole in a field precesses and radiates continuously.  The total energy radiated depends on the angle $\theta$ between spin and field: maximum if $\theta=\pi$ (anti-aligned), zero if $\theta=0$ (aligned), and intermediate values for intermediate angles.  The radiation frequency is always $\omega$.

\noindent\textbf{Quantum picture:}
\begin{itemize}[nosep]
    \item A spin in $\ket{\uparrow}$ (upper level) can emit a photon of energy $\omega$ and transition to $\ket{\downarrow}$.
    \item A spin in $\ket{\downarrow}$ cannot emit (no lower state available).
    \item A spin tilted at an angle (e.g.\ along $x$, in state $\frac{1}{\sqrt{2}}(\ket{\uparrow}+\ket{\downarrow})$) does \emph{not} emit a photon of half-energy $\omega/2$.  Instead: with probability $|\alpha|^2=\frac{1}{2}$ it emits a photon of full energy $\omega$, and with probability $|\beta|^2=\frac{1}{2}$ it emits nothing.
    \item The \textbf{average} energy emitted equals the classical expectation ($\frac{\omega}{2}$ for a horizontal spin)---but each individual event is all-or-nothing.
\end{itemize}

% ============================================================
\section{2. Collapse of the Wave Function}
% ============================================================

Between measurements, the state evolves unitarily via the Schr\"odinger equation.  At the moment of measurement of observable $L$:

\begin{enumerate}[nosep]
    \item The system is in some state $\ket{\psi}=\sum_i \alpha_i\ket{\lambda_i}$ (expanded in eigenstates of $L$).
    \item Measurement yields eigenvalue $\lambda_k$ with probability $|\alpha_k|^2$.
    \item Immediately after, the state \textbf{collapses} to $\ket{\lambda_k}$.  All other terms in the superposition become irrelevant---only one eigenstate ``survives.''
\end{enumerate}

This ``collapse'' is \emph{not} governed by the Schr\"odinger equation.  It appears to be a second, separate rule for how states change.  This is unsatisfying: apparatuses are themselves quantum systems.  Why should there be two rules---one (Schr\"odinger) for undisturbed evolution, another (collapse) for measurement?

The resolution: if we treat the apparatus as a quantum system too, then the combined system (system + apparatus) evolves by the Schr\"odinger equation alone.  Collapse emerges naturally from analyzing the composite system.  This motivates the study of \textbf{composite quantum systems}.

% ============================================================
\section{3. Combining Quantum Systems: The Tensor Product}
% ============================================================

\begin{definition}[Tensor product]
Given system $\mathcal{A}$ with state space $\mathcal{S}_A$ (dimension $d_A$, basis $\{\ket{a}\}$) and system $\mathcal{I}$ with state space $\mathcal{S}_I$ (dimension $d_I$, basis $\{\ket{i}\}$), the \textbf{combined} state space is the tensor product
\[
    \mathcal{S}_{AI} = \mathcal{S}_A \otimes \mathcal{S}_I,
\]
with dimension $d_A \times d_I$.  Its basis vectors are all pairs $\ket{a,i}\equiv\ket{a}\otimes\ket{i}$.
\end{definition}

\subsection{3.1 Inner product on the tensor product}
\[
    \braket{b,j}{a,i} = \delta_{ab}\,\delta_{ij}.
\]
Two composite basis states are orthogonal unless \emph{both} labels match.

\subsection{3.2 General state}
\[
    \ket{\psi} = \sum_{a,i} \alpha_{ai}\,\ket{a,i}, \qquad \alpha_{ai}\in\mathbb{C}.
\]

% ============================================================
\section{4. Two Spins}
% ============================================================

Label the spins $\sigma$ (Alice) and $\tau$ (Bob).  Each lives in $\mathbb{C}^2$.  The combined space is $\mathbb{C}^2\otimes\mathbb{C}^2 = \mathbb{C}^4$ with basis:
\[
    \ket{\uparrow\uparrow},\quad \ket{\uparrow\downarrow},\quad \ket{\downarrow\uparrow},\quad \ket{\downarrow\downarrow},
\]
where the first arrow refers to $\sigma$ and the second to $\tau$.

A general two-spin state:
\[
    \ket{\psi} = \alpha_{\uparrow\uparrow}\ket{\uparrow\uparrow}
    + \alpha_{\uparrow\downarrow}\ket{\uparrow\downarrow}
    + \alpha_{\downarrow\uparrow}\ket{\downarrow\uparrow}
    + \alpha_{\downarrow\downarrow}\ket{\downarrow\downarrow}.
\]

% ============================================================
\section{5. Product States vs.\ Entangled States}
% ============================================================

\begin{definition}[Product state]
A state $\ket{\psi}$ of the composite system is a \textbf{product state} if it can be written as
\[
    \ket{\psi} = \ket{\phi}_\sigma \otimes \ket{\chi}_\tau,
\]
equivalently, the coefficients factorize: $\alpha_{ab} = \phi_a\,\chi_b$.
\end{definition}

\begin{definition}[Entangled state]
A state that is \textbf{not} a product state is called \textbf{entangled}.
\end{definition}

\subsection{5.1 Parameter counting}
A general two-spin state has 4 complex coefficients $= 8$ real parameters.  Subtract 1 for normalization and 1 for overall phase: \textbf{6 real parameters}.

A product state is specified by one single-spin state (2 real params) for $\sigma$ and one (2 real params) for $\tau$: \textbf{4 real parameters}.

Since $6>4$, the product states form a \emph{lower-dimensional} subset: \textbf{most states are entangled}.

\subsection{5.2 Example: product state}
$\ket{\uparrow}\otimes\ket{\rightarrow} = \frac{1}{\sqrt{2}}\bigl(\ket{\uparrow\uparrow}+\ket{\uparrow\downarrow}\bigr).$  
Alice's spin is definitely up; Bob's is in state $\ket{\rightarrow}$.  Each subsystem has its own well-defined state.

\subsection{5.3 Example: entangled state (singlet)}
\[
    \ket{\text{singlet}} = \frac{1}{\sqrt{2}}\bigl(\ket{\uparrow\downarrow}-\ket{\downarrow\uparrow}\bigr).
\]
This cannot be written as a product.  Neither spin has a definite individual state.

\noindent\textbf{Remarkable property:} The singlet looks the same in \emph{every} basis.  If you rewrite $\ket{\uparrow},\ket{\downarrow}$ in terms of $\ket{\rightarrow},\ket{\leftarrow}$, you get $\frac{1}{\sqrt{2}}(\ket{\rightarrow\leftarrow}-\ket{\leftarrow\rightarrow})$.  In terms of in/out: $\frac{1}{\sqrt{2}}(\ket{\text{in},\text{out}}-\ket{\text{out},\text{in}})$.  This basis-independence is a hallmark of maximal entanglement.

% ============================================================
\section{6. Observables on Composite Systems}
% ============================================================

\subsection{6.1 Alice's observable in the combined space}
An observable $L$ belonging to Alice acts on Alice's indices and is the identity on Bob's:
\[
    L_{a'b',\,ab} = L_{a'a}\,\delta_{b'b}.
\]

\subsection{6.2 Expectation values in a product state}
In a product state $\ket{\psi}=\ket{\phi}_A\otimes\ket{\chi}_B$:
\[
    \expval{L}_\psi = \expval{L}_\phi,
\]
i.e.\ Alice's expectation values are the same as if Bob did not exist.  Likewise for Bob's observable $M$:
\[
    \expval{M}_\psi = \expval{M}_\chi.
\]

\subsection{6.3 Product of observables}
For a product state:
\[
    \expval{LM} = \expval{L}\,\expval{M}.
\]
The expectation value of the product equals the product of expectation values---the variables are \textbf{uncorrelated}.

% ============================================================
\section{7. Correlations and Entanglement}
% ============================================================

\begin{definition}[Correlation]
The \textbf{correlation} between Alice's observable $L$ and Bob's observable $M$ is
\[
    C(L,M) = \expval{LM} - \expval{L}\expval{M}.
\]
\end{definition}

\begin{itemize}[nosep]
    \item For product states: $C(L,M)=0$ for all $L,M$.
    \item For entangled states: there exists at least one pair $(L,M)$ with $C(L,M)\neq 0$.
\end{itemize}

\subsection{7.1 All individual expectation values vanish}
In the singlet state, compute $\expval{\sigma_z}$:
\[
    \expval{\sigma_z} = \tfrac{1}{2}\bigl(\bra{\uparrow\downarrow}-\bra{\downarrow\uparrow}\bigr)\sigma_z\bigl(\ket{\uparrow\downarrow}-\ket{\downarrow\uparrow}\bigr).
\]
$\sigma_z$ acts only on the first index: $\sigma_z\ket{\uparrow\downarrow}=+\ket{\uparrow\downarrow}$, $\sigma_z\ket{\downarrow\uparrow}=-\ket{\downarrow\uparrow}$.  So:
\[
    \expval{\sigma_z} = \tfrac{1}{2}\bigl(\braket{\uparrow\downarrow}{\uparrow\downarrow}+\braket{\downarrow\uparrow}{\downarrow\uparrow}\bigr) = \tfrac{1}{2}(1-1)... 
\]
Working through cross terms: the result is $\expval{\sigma_z}=0$.  Similarly $\expval{\sigma_x}=0$ and $\expval{\sigma_y}=0$ (since $\sigma_x$ and $\sigma_y$ flip the first index, producing states orthogonal to the originals).

But for a \emph{single} spin in \emph{any} state, there always exists a direction $\hat{n}$ with $\expval{\vec{\sigma}\cdot\hat{n}}=+1$.  The fact that all components vanish simultaneously proves the singlet \textbf{cannot be a product state}.

\subsection{7.2 Singlet state correlations}
\begin{align}
    \expval{\sigma_z} &= 0, & \expval{\tau_z} &= 0, \\
    \expval{\sigma_z\,\tau_z} &= -1. &&
\end{align}
\textbf{Calculation of $\expval{\sigma_z\tau_z}$}: $\tau_z$ acts on the second index ($+1$ for $\uparrow$, $-1$ for $\downarrow$), then $\sigma_z$ acts on the first.  On $\ket{\uparrow\downarrow}$: $\sigma_z\tau_z = (+1)(-1)=-1$.  On $\ket{\downarrow\uparrow}$: $\sigma_z\tau_z = (-1)(+1)=-1$.  Both terms give $-1$, so $\expval{\sigma_z\tau_z}=-1$.  In fact, the singlet is an \emph{eigenstate} of $\sigma_z\tau_z$ with eigenvalue $-1$.

\medskip
\noindent Correlation: $C(\sigma_z,\tau_z) = -1 - (0)(0) = -1 \neq 0$.  Similarly, $\expval{\sigma_x\tau_x}=-1$ and $\expval{\sigma_y\tau_y}=-1$.  But cross-products like $\expval{\sigma_x\tau_z}=0$.

\medskip
\noindent\textbf{Physical meaning:} If you measure the \emph{same} component of both spins, they are always opposite.  If you measure \emph{different} components, there is no correlation.

% ============================================================
\section{8. Summary}
% ============================================================

\begin{enumerate}[nosep]
    \item Photon emission is all-or-nothing: a tilted spin emits a full-energy photon with some probability, or nothing.  The average matches the classical prediction.
    \item Wave function collapse is not a separate law---it emerges when the apparatus is included as part of the quantum system.
    \item Composite systems are described by the \textbf{tensor product} of individual state spaces ($\dim = d_A \times d_I$).
    \item A \textbf{product state} factorizes (4 real params for two spins); an \textbf{entangled state} does not (6 real params).  Most states are entangled.
    \item In a product state, subsystem observables are uncorrelated: $\expval{LM}=\expval{L}\expval{M}$.
    \item Nonzero correlations $C(L,M)\neq 0$ signal entanglement.
    \item The singlet is maximally entangled: all individual $\expval{\sigma_i}=0$, but $\expval{\sigma_i\tau_i}=-1$ for $i=x,y,z$.  It looks the same in every basis.
    \item Entanglement correlations resemble classical correlations (penny/dime analogy) but cannot be simulated by classical computers without nonlocal wiring---a preview of Bell's theorem.
\end{enumerate}

\end{document}
